%==============================================================================
% Sjabloon poster bachproef
%==============================================================================
% Gebaseerd op document class `a0poster' door Gerlinde Kettl en Matthias Weiser
% Aangepast voor gebruik aan HOGENT door Jens Buysse en Bert Van Vreckem

\documentclass[a0,portrait]{hogent-poster}

% Info over de opleiding
\course{Bachelorproef}
\studyprogramme{toegepaste informatica}
\academicyear{2025-2026}
\institution{Hogeschool Gent, Valentin Vaerwyckweg 1, 9000 Gent}

% Info over de bachelorproef
\title{Hoe kan een Privileged Access Management tool bijdragen aan de beveiliging van servers}
\subtitle{}
\author{Brent Lissens}
\email{brent.lissens@student.hogent.be}
\supervisor{Lena De Mol \& Nathalie Declercq}
\cosupervisor{Lien Gillis (AE)}

% Indien ingevuld, wordt deze informatie toegevoegd aan het einde van de
% abstract. Zet in commentaar als je dit niet wilt.
\specialisation{Systeem- en Netwerkbeheer}
\keywords{Privileged Access Management, PAM, Serverbeveiliging, Zero Trust, NIS2}
\projectrepo{https://github.com/user/repo}

\begin{document}

\maketitle

\begin{abstract}
	Privileged accounts vormen een aantrekkelijk doelwit voor cyberaanvallers omdat zij uitgebreide rechten hebben op kritieke systemen. Het beveiligen van deze accounts is daarom essentieel voor organisaties die hun servers en infrastructuur willen beschermen.
	
	In deze bachelorproef werd onderzocht hoe een Privileged Access Management (PAM)-oplossing kan bijdragen aan de beveiliging van servers. Hiervoor werd een vergelijkende studie uitgevoerd van CyberArk, BeyondTrust, ManageEngine PAM360 en Securden. Daarnaast werd een proof-of-concept gerealiseerd met PAM360 en Securden om de oplossingen in een praktijkomgeving te evalueren.
\end{abstract}

\begin{multicols}{2} % This is how many columns your poster will be broken into, a portrait poster is generally split into 2 columns

\section{Introductie}

Privileged accounts beschikken over uitgebreide rechten op systemen en vormen daardoor een belangrijk doelwit voor cyberaanvallers. Traditioneel worden deze accounts vaak beheerd met statische wachtwoorden en beperkte monitoring, wat het risico op misbruik verhoogt.

Deze bachelorproef onderzoekt hoe een Privileged Access Management-oplossing kan bijdragen aan de beveiliging van servers. Daarbij wordt nagegaan welke functionaliteiten PAM-oplossingen aanbieden en hoe deze organisaties kunnen ondersteunen bij het beschermen van kritieke systemen en het naleven van beveiligingsrichtlijnen zoals NIS2.

Belangrijke PAM-functionaliteiten zijn onder andere credential vaulting, session management, just-in-time access, toegangscontrole en auditing.

\section{Experimenten}

Het onderzoek bestond uit twee onderdelen.

Enerzijds werd een vergelijkende studie uitgevoerd van vier PAM-oplossingen:
\begin{itemize}
	\item CyberArk
	\item BeyondTrust
	\item ManageEngine PAM360
	\item Securden Unified PAM
\end{itemize}

Deze oplossingen werden geëvalueerd op basis van criteria zoals credential vaulting, session management, just-in-time access, toegangscontrole, Active Directory-integratie en implementatiecomplexiteit.

Anderzijds werd een proof-of-concept opgezet waarbij PAM360 en Securden werden geïmplementeerd in een gecontroleerde testomgeving. Hiervoor werd een gesegmenteerd netwerk ingericht met Windows- en Linux-servers, Active Directory en een Fortinet-firewall.

Tijdens de proof-of-concept werden onder andere SSH- en RDP-toegang, credential vaulting, sessieregistratie en rapportering getest.


\section{Conclusies}

Uit de vergelijkende studie blijkt dat alle onderzochte oplossingen bijdragen aan het beveiligen van privileged accounts, maar dat er duidelijke verschillen bestaan op vlak van functionaliteit, gebruiksgemak en implementatiecomplexiteit.

CyberArk en BeyondTrust bieden de meest uitgebreide functionaliteiten en zijn geschikt voor complexe enterprise-omgevingen. PAM360 combineert een brede functionaliteit met een relatief toegankelijke implementatie. Securden onderscheidt zich door zijn eenvoud en gebruiksvriendelijke aanpak.

De proof-of-concept bevestigde dat PAM-oplossingen een belangrijke meerwaarde bieden door wachtwoorden centraal te beheren, privileged sessies te monitoren en toegang strikt te controleren.

\section{Toekomstig onderzoek}

Verder onderzoek kan zich richten op:

\begin{itemize}
	\item PAM binnen cloud- en hybride omgevingen;
	\item geavanceerde Just-in-Time toegangstechnieken;
	\item integratie met Zero Trust architecturen;
	\item verdere automatisering van privileged access governance;
	\item kosten-batenanalyse van PAM-implementaties op lange termijn.
\end{itemize}

\end{multicols}
\end{document}